\documentclass[../../../include/open-logic-section]{subfiles}

\begin{document}

\olfileid[fr]{sth}{story}{rus}
\olsection{Le paradoxe de Russell (à nouveau)}

Dans la \olref[sfr][][]{part}, nous avons travaillé avec une théorie
naïve des ensembles. Mais, selon une conception \emph{très}
naïve, les ensembles ne sont que les extensions des prédicats. Cette
idée naïve imposerait le principe suivant :

\begin{defish}
  \emph{Compréhension naïve.} $\Setabs{x}{\phi(x)}$ existe pour toute formule $\phi$.
\end{defish}

Si séduisant que soit ce principe, on peut démontrer qu'il est
contradictoire.
Nous l'avons vu à la \olref[sfr][set][rus]{sec}, mais ce résultat est
si important et si simple qu'il vaut la peine de le répéter. Mot pour
mot.

\begin{thm}[Paradoxe de Russell]
Il n'existe pas d'ensemble $R = \Setabs{x}{x \notin x}$.
\end{thm}

\begin{proof}
Si $R = \Setabs{x}{x \notin x}$ existe, alors $R \in R$ si et
seulement si $R \notin R$, ce qui est contradictoire.
\end{proof}

Russell découvrit ce résultat en juin 1901. (Il ne formula cependant
pas le paradoxe tout à fait sous la forme que nous venons de présenter,
puisqu'il examinait la théorie des ensembles de Frege, telle qu'elle
est exposée dans les \emph{Grundgesetze}. Nous y reviendrons à la
\olref[blv]{sec}.) Russell écrivit à Frege le 16 juin 1902 pour lui
expliquer la contradiction dans son système. Pour leur correspondance
et quelques éléments de contexte, voir \citet[pp.~124--8]{Heijenoort1967}.

Soulignons que cette démonstration de deux lignes relève de la
\emph{logique pure}. Certes, notre notation $\Setabs{x}{x \notin x}$
fait implicitement appel à un axiome (non logique ?)\ :
l'extensionnalité. En effet, $\Setabs{x}{\phi(x)}$ désigne
\emph{l'unique} ensemble (par extensionnalité) des objets qui
satisfont~$\phi$, s'il en existe un. Mais nous pouvons éviter jusqu'à
cette trace d'extensionnalité en énonçant simplement le résultat ainsi :
\emph{il n'existe aucun ensemble dont les membres soient exactement
les ensembles qui ne s'appartiennent pas à eux-mêmes}. Cela n'a
d'ailleurs pas grand-chose à voir avec les ensembles. Comme Russell
l'observait lui-même, le même raisonnement conduit à conclure :
\emph{aucun homme ne rase exactement les hommes qui ne se rasent pas
eux-mêmes}. Ou encore : \emph{aucun carlin ne renifle exactement les
carlins qui ne se reniflent pas eux-mêmes}. Et ainsi de suite. Sous
une forme schématique, le résultat est simplement :
\[
\lnot \exists x \forall z(Rzx \liff \lnot R zz).
\]
Il s'agit là d'un théorème (schéma) de la logique du premier ordre.
Par conséquent, modifier seulement notre théorie des ensembles ne
suffirait pas à éviter le paradoxe de Russell : il surgit avant même
que nous \emph{n'abordions} la théorie des ensembles. Si nous
utilisons la logique du premier ordre (classique), nous devons
simplement \emph{accepter} qu'il n'existe pas d'ensemble
$R = \Setabs{x}{x\notin x}$.

Il en résulte ceci. Si l'on veut admettre la compréhension naïve tout
en \emph{évitant} la contradiction, on ne peut pas se contenter de
modifier la \emph{théorie des ensembles}. Il faudrait plutôt remanier
la \emph{logique} elle-même.

Bien sûr, des théories des ensembles fondées sur des logiques non
classiques ont été proposées. Mais elles sont, pour le moins, non
standard. L'approche habituelle du paradoxe de Russell consiste à le
traiter comme une simple démonstration de non-existence, puis à
chercher comment en tenir compte. C'est l'approche que nous suivrons.

\end{document}
